\chapter{Literature Review}

\section{Domain Background}

Digital heritage systems require high recognition reliability under uncontrolled
acquisition conditions. Museum deployments add constraints including privacy,
latency, multilingual support, and explainable outcomes.

Unlike lab datasets, museum environments introduce dynamic human behavior: people
move in front of displays, lighting varies across halls, glass reflections appear
unexpectedly, and objects are photographed from oblique angles. A model that performs
well in clean, centered-image conditions can fail abruptly in these scenarios. The
literature review is therefore interpreted through deployment realism rather than
isolated benchmark ranking.

\section{State of the Art in Software Solutions}

Existing solutions include cloud classification APIs, AR image-retrieval systems, and
on-device CNN classifiers. Each approach trades off latency, privacy, and robustness.
Marker-based AR remains common in heritage apps but is sensitive to viewing angle and
lighting and scales poorly to large artifact sets, as documented in case studies of
Egyptian cultural-heritage AR~\citep{abdelhady2025enhancing}.

\section{Survey of AI Approaches}

\subsection{On-Device Object Detection}
Single-stage detectors of the YOLO family provide a strong speed/accuracy balance for
mobile deployment, and export cleanly to TensorFlow Lite for on-device inference with
end-to-end non-maximum suppression baked into the graph~\citep{yolov8}. Comparative
studies of museum-exhibit recognition from video-derived datasets report that YOLOv8
reaches very high mean Average Precision while remaining light enough for real-time
mobile use, outperforming heavier backbones such as VGG-16 and ResNet variants on the
speed/accuracy trade-off~\citep{ipalakova2025comparative}. Adaptive CNN approaches
have likewise been used to enrich museum interaction through automatic artifact
recognition~\citep{wen2024enhancing}.

\subsection{Foundation-Model Auto-Annotation}
Manually annotating tens of thousands of frames is prohibitive for a student-led
project, motivating \emph{auto-annotation} with open-vocabulary, text-prompted
detectors such as YOLO-World~\citep{yoloworld}. Such foundation models handle generic
concepts well but degrade on fine-grained, domain-specific objects---precisely the
near-identical statues that dominate the GEM collection---producing systematically
mislocalized or sibling-confused labels. This limitation is the central motivation for
the data-quality study in Chapter~5.

\subsection{Retrieval-Augmented Generation and On-Device LLMs}
Retrieval-Augmented Generation (RAG) grounds a language model's output in passages
retrieved from a knowledge base rather than its parametric memory, reducing
fabrication in knowledge-intensive tasks~\citep{lewis2020rag}. Recent work also brings
context-aware LLM assistants into AR settings~\citep{qorbani2025teaching}. On the
serving side, quantized models run locally through runtimes such as
Ollama~\citep{ollama}, while multilingual sentence embeddings make a single bilingual
(English/Arabic) retrieval index practical. Robustness challenges specific to heritage
capture---low contrast, reduced lighting for pigment preservation, and crowd
occlusion---are active research topics~\citep{liu2025yolov8,dai2024occlusion}.

\section{Identification of Research Gaps}

The review identifies three combined gaps that this thesis addresses: (i) label
quality for fine-grained, domain-specific detection; (ii) honest evaluation of
video-frame datasets, where random splits leak near-duplicate frames; and (iii)
grounded, bilingual generation for public-facing museum guides. Table~\ref{tab:gaps}
maps prior approaches to how \TimeLens{} fills each gap.

\begin{table}[htbp]
\centering
\caption{Synthesis of prior work and gap fit.}
\label{tab:gaps}
\begin{tabularx}{\textwidth}{@{}>{\raggedright\arraybackslash}p{0.24\textwidth}>{\raggedright\arraybackslash}p{0.20\textwidth}>{\raggedright\arraybackslash}p{0.22\textwidth}>{\raggedright\arraybackslash}X@{}}
\toprule
\textbf{Approach} & \textbf{Strength} & \textbf{Limitation} & \textbf{How TimeLens addresses it}\\
\midrule
Cloud classification APIs & Easy integration, strong models & Latency, privacy, connectivity dependence & On-device YOLOv8n detection (no per-frame round-trip)\\
Foundation-model auto-annotation (YOLO-World) & Cheap labels at scale & Fails on fine-grained, domain-specific artifacts & Label-quality iteration: spatial cleaning $+$ hand annotation\\
Random-split evaluation on video frames & Simple to run & Frame leakage inflates metrics & Honest disclosure $+$ conservative held-out; future video-level split\\
Ungrounded LLM guides & Fluent, engaging & Hallucinate facts; weak Arabic & Bilingual RAG grounding (ChromaDB $+$ Gemma 4 E2B)\\
\bottomrule
\end{tabularx}
\end{table}
